% ch06-censoring --- untrusted and censored deal times: error models,
% right censoring under history dependence, interval censoring, and the
% mean-intensity (MBPP) route of Rizoiu et al. (2022). Owns defect (C4).

\section{What the dates are, and what we choose to trust}\label{ch06:sec:defect}

\Cref{ch:data} recorded, as defect (C4), that the event times $t_m$ of
\cref{def:ch02:market} are not delivered to us. What PitchBook delivers is a
recorded date $\hat t_m$ produced by an editorial process: announced versus
closed dates that differ by days, back-filled records whose dates were
reconstructed after the fact (often rounded to a month or quarter boundary),
and duplicate rows collapsed by the hygiene rules of
\cref{prot:ch02:pit}. Every continuous-time likelihood in
\cref{ch:poisson,ch:hawkes} consumes $t_m$ as if it were exact; this chapter
decides what to do when it is not. The nonparametric intensity models of
\cref{ch:nonparam} are valid only on trusted, uncensored dates and have no
interval-censored counterpart here; when dates cannot be trusted, that
regime must fall back to the binned models of this chapter. Our thesis,
defended throughout: \emph{the correct response to untrusted positions is to choose
the observation resolution honestly, and to use likelihoods that consume
only what we trust --- counts of events on intervals, never positions within
them.} The count is an editorial fact we can audit; the timestamp is an
editorial choice we cannot.

We formalize the corruption with three assumptions, kept deliberately weak.

\begin{assumption}[Bounded date jitter]\label{asm:ch06:jitter}
Each deal $m$ has a true time $t_m$ and a recorded time
$\hat t_m = t_m + \eps_m$ with $\abs{\eps_m} \le \delta_m \le \delta$ for a
known trust radius $\delta$. Nothing else is assumed about the errors
$\eps_m$: they may be dependent on each other, on the process, and on the
marks; they need not be centered. The bound $\delta$ is an output of the
date audit (\cref{ch06:sec:decisions}), not a modeling choice.
\end{assumption}

\begin{assumption}[Recording conventions]\label{asm:ch06:convention}
The recorded date decomposes as
$\hat t_m = t_m + b(r_m) + \eps_m$, where $r_m \in
\{\text{closed},\allowbreak \text{announced},\allowbreak
\text{backfilled}\}$ is an observable record type, $b(\cdot)$ is a
systematic convention offset with $b(\text{closed}) = 0$, and backfilled
records may additionally exhibit \emph{heaping}: rounding to calendar
boundaries, visible as excess mass on month and quarter starts. After
auditing to a single convention (subtracting an estimated $b$, flagging
heaped records), the residual error satisfies \cref{asm:ch06:jitter} with a
type-specific $\delta_m$.
\end{assumption}

\begin{assumption}[Count completeness]\label{asm:ch06:complete}
After the deduplication of \cref{prot:ch02:pit}, every in-window deal of the
tracked process appears exactly once in the record. Corruption lives in the
\emph{position} $\hat t_m$, never in the \emph{multiplicity}. (Thinning of
the record --- events missing entirely --- is a different defect: (C6) for
out-of-universe attribution, (C3) for unobserved exits. This chapter is
orthogonal to both; see \cref{ch06:sec:decisions}.)
\end{assumption}

The deliberate weakness of \cref{asm:ch06:jitter} --- adversarial, possibly
dependent errors --- is what makes the induced modeling choice sharp. If the
$\eps_m$ were i.i.d.\ with a known density we could deconvolve; we refuse to
assume that, because announced-versus-closed offsets are systematic, not
noise. Under adversarial bounded jitter the only event-level statements that
survive are set-membership statements, and the following proposition is the
complete list of what we keep.

\begin{proposition}[What bounded jitter leaves invariant]\label{prop:ch06:jitterrobust}
Let $0 = o_0 < o_1 < \cdots < o_J = \Tend$ be a partition into bins
$I_j = (o_{j-1}, o_j]$ of width $\horizon = o_j - o_{j-1}$, let
$C_j = \#\{m : t_m \in I_j\}$ be the true bin counts and
$\hat C_j = \#\{m : \hat t_m \in I_j\}$ the observed ones. Under
\cref{asm:ch06:jitter,asm:ch06:complete}:
\begin{enumerate}[label=(\roman*)]
\item every event observed in bin $j$ has its true time in the dilated bin
$(o_{j-1}-\delta,\, o_j+\delta]$;
\item $\abs{\hat C_j - C_j} \le
\#\{m : \min_k \abs{t_m - o_k} \le \delta\}$ restricted to the two
boundaries of $I_j$: only events within $\delta$ of a bin boundary can be
miscounted, and only into an adjacent bin;
\item if the process has locally bounded mean rate
$\sup_t \E[\lam(t)] \le \bar\lam_{\sup}$, then
$\E\abs{\hat C_j - C_j} \le 4\delta\,\bar\lam_{\sup}$, so the expected
miscounted \emph{fraction} per bin is of order $4\delta/\horizon$;
\item counts over unions of consecutive bins are more robust still:
transfers across interior boundaries cancel, leaving only the two outer
boundaries.
\end{enumerate}
\end{proposition}

\begin{proof}
(i) is immediate from $\abs{\hat t_m - t_m}\le\delta$. (ii): an event with
$t_m \in I_j$ and $\hat t_m \notin I_j$ requires $\hat t_m$ on the far side
of $o_{j-1}$ or $o_j$, hence $\abs{t_m - o_k}\le\delta$ for a boundary of
$I_j$; conversely for events entering $I_j$. (iii): by (ii) and the mean-measure
bound $\E\,\#\{m: t_m \in A\} = \int_A \E[\lam] \le
\abs{A}\,\bar\lam_{\sup}$, the expected number of events within $\delta$ of
a fixed boundary is at most $2\delta\,\bar\lam_{\sup}$; each bin has two
boundaries. Comparing with $\E C_j \asymp \horizon\bar\lam$ gives
the fraction. (iv): an event transferred from $I_j$ to $I_{j+1}$ leaves the
count of $I_j \cup I_{j+1}$ unchanged.
\end{proof}

The induced choice is now forced. Any inference that reads event
\emph{positions} at scales comparable to $\delta$ is reading editorial
noise; any inference that reads \emph{counts} on bins of width
$\horizon \gg \delta$ is reading the process, up to a boundary-transfer
error of controlled order $4\delta/\horizon$ that vanishes as the resolution
coarsens. We therefore commit to modeling at resolution no finer than the
trust radius --- in practice $\horizon \ge$ several multiples of $\delta$ ---
at which point \emph{the data is, by construction, interval-censored count
data}, and the rest of this chapter builds the likelihoods that consume
exactly that and nothing more. One reassurance: the
mark factor of \cref{eq:ch01:groundmark} conditions on a deal occurring and
compares candidates within a risk set, so it consumes dates only through
risk-set membership and coarse ordering; jitter at the scale of days is
invisible to it (\cref{ch:ranking} quantifies this). It is the ground
process --- the timing factor --- that absorbs (C4), which is why this
chapter is about intensities.

\section{Right censoring with history-dependent intensities}\label{ch06:sec:rightcens}

Before interval censoring, we settle right censoring --- defect (C1) --- for
the Hawkes case, because a doubt naturally arises there that does not arise
for Poisson models. For an NHPP the censored-window contribution
$\exp(-\int\lam)$ involves a deterministic function, and the argument of
\cref{ch:poisson} is elementary. For a Hawkes process the intensity is
random and feeds on the history: one might worry that the survival factor
for the window $(t_{\text{last}}, \Tend]$ secretly depends on events we have
not seen, or that the ABGK independent-censoring argument breaks. It does
not, and the reason --- predictability --- is worth stating precisely,
because it is the same property that will \emph{fail} to rescue us under
interval censoring.

\begin{assumption}[Independent, noninformative censoring]\label{asm:ch06:indcens}
Let $N$ be a counting process with $\Ft$-intensity
$\atrisk(t)\lam(t;\thetamodel)$ as in \cref{def:ch02:intensity}. There is an
enlarged filtration $\mathcal{G}_t \supseteq \Ft$ carrying the censoring
information such that (i) the $\mathcal{G}_t$-intensity of $N$ is still
$\atrisk(t)\lam(t;\thetamodel)$; (ii) the censoring indicator
$C(t)\in\{0,1\}$ (e.g.\ $C(t)=\ind{t \le \tau}$ for a censoring time
$\tau$) is $\mathcal{G}_t$-predictable; (iii) the censoring mechanism's
factor in the full likelihood does not involve $\thetamodel$
\citep[III.2]{abgk1993}. Administrative censoring at the fixed date
$\Tend$ (June 2024) satisfies all three with $\mathcal{G}=\F$ and
$C(t)=\ind{t\le\Tend}$.
\end{assumption}

\begin{theorem}[Censored-window validity for history-dependent intensities]\label{thm:ch06:rightcensor}
Under \cref{asm:ch06:indcens}, with $\lam$ any $\mathcal{G}$-predictable
intensity --- in particular the Hawkes intensity
$\lam(t) = \mub(t) + \sum_{t_m < t} \kernel(t - t_m)$, which depends on the
realized history:
\begin{enumerate}[label=(\roman*)]
\item the observed process $N^{o}(t) = \int_0^t C(u)\dd N(u)$ has
$\mathcal{G}_t$-intensity $C(t)\,\atrisk(t)\,\lam(t;\thetamodel)$;
\item the likelihood of the observed data factorizes as the Jacod
likelihood \cref{eq:ch02:jacod} evaluated with intensity
$C\atrisk\lam$, times a factor free of $\thetamodel$; the first factor is
therefore a genuine partial likelihood, and a censored spell contributes
exactly its survival factor
$\exp\!\big({-}\!\int_{\text{spell}} \lam\big)$ --- no term dropped, no
term invented;
\item the score of this partial likelihood is a $\mathcal{G}$-martingale
evaluated at $\Tend$, so the estimation theory of \cref{ch:foundations}
(unbiased score, information identities, asymptotic normality) applies
unchanged.
\end{enumerate}
\end{theorem}

\begin{proof}
(i) $M = N - \int \atrisk\lam\dd u$ is a local $\mathcal{G}$-martingale by
\cref{asm:ch06:indcens}(i). The process $C$ is bounded and
$\mathcal{G}$-predictable by (ii), so the stochastic integral
$\int_0^t C \dd M = N^{o}(t) - \int_0^t C\atrisk\lam \dd u$ is again a local
$\mathcal{G}$-martingale; since $\int C\atrisk\lam\dd u$ is predictable and
nondecreasing, it is the compensator of $N^{o}$, i.e.\ $C\atrisk\lam$ is the
intensity of the observed process. Observe what the proof used: boundedness
and predictability of $C$, predictability of $\lam$ --- \emph{nowhere} that
$\lam$ is deterministic. History dependence is immaterial because
predictability is preserved under enlargement by
\cref{asm:ch06:indcens}(i); this is the ABGK III.2 argument verbatim.
(ii) is the standard likelihood factorization for counting processes with
the intensity from (i) \citep[II.7, III.2]{abgk1993}: the full likelihood is
the product integral of the observed-jump and survival factors for $N^{o}$
times the censoring factor, and (iii) removes the latter from the
maximization. (iii) follows because the score is
$\int_0^{\Tend} \partial_{\thetamodel}\log\lam \,(\dd N^{o} -
C\atrisk\lam\dd u)$, a stochastic integral of a predictable process against
the martingale of (i).
\end{proof}

What, then, is genuinely different from the Poisson case? Only this: the
survival factor is now \emph{random} --- but it is a functional of the
\emph{observed} data alone. The next proposition makes the spec's
``depends on the unobserved post-censoring events only through ---
nothing'' precise.

\begin{proposition}[Locality of the compensator]\label{prop:ch06:locality}
Let $\tau \le \Tend$ be the censoring time of \cref{asm:ch06:indcens} and
let $\lam$ be the Hawkes intensity with causal kernel ($\kernel$ supported
on $(0,\infty)$). For every $t \le \tau$, $\lam(t)$ is measurable with
respect to $\Hist_{t^-}$, the history of events \emph{strictly before} $t$;
consequently the integrated compensator
$\Lam(\tau) = \int_0^{\tau}\lam(u)\dd u$ is measurable with respect to the
observed history $\Hist_{\tau}$, and the survival factor
$e^{-\Lam(\tau)}$ is a deterministic functional of the data we possess.
Events in $(\tau, \infty)$ --- observed or not --- enter neither the event
terms nor the compensator of \cref{thm:ch06:rightcensor}(ii).
\end{proposition}

\begin{proof}
$\lam(t) = \mub(t) + \sum_{t_m < t}\kernel(t - t_m)$ involves only event
times $t_m < t \le \tau$, all of which are uncensored under
\cref{asm:ch06:complete}; measurability of the integral follows from joint
measurability of the predictable path $u \mapsto \lam(u)$ on $(0,\tau]$ and
Tonelli. The claim is exactly the statement that the compensator up to a
stopping time is determined by the pre-stopping-time path --- the defining
property of predictability, which the Hawkes feedback respects because the
kernel is causal.
\end{proof}

\begin{warning}[Completeness of the pre-censoring record is now load-bearing]\label{warn:ch06:missingpast}
For a Poisson model, an event missing from the record costs exactly its own
two terms (one jump term, one sliver of compensator); the rest of the
likelihood is untouched. For a Hawkes model, a missing past event corrupts
the intensity \emph{path} at every later time: the reconstructed
$\lam(t)$ omits a kernel contribution $\kernel(t - t_m)$, so every
subsequent event term and the entire survival factor are evaluated at the
wrong intensity. The bias has no universal sign --- the fitted baseline
$\mub$ and branching scale $\branch$ reapportion the missing excitation ---
which is precisely why \cref{asm:ch06:complete} (counts complete, positions
untrusted) is the boundary of this chapter's mandate: it protects the
history that the intensity feeds on. Where completeness itself fails, we
are in (C3)/(C6) territory and no date discipline repairs it.
\end{warning}

\section{Interval censoring for Poisson models: exactness}\label{ch06:sec:icpoisson}

Coarsening event times to bin counts is interval censoring by construction.
For the Poisson rung of the ladder this costs \emph{no approximation at
all} --- only resolution. We restate the binned-likelihood theorem of
\cref{ch:poisson} in the form needed here, and prove it by a route chosen
deliberately: the same configuration integral, evaluated for a
history-dependent intensity in the next section, is exactly where
tractability fails.

\begin{proposition}[Binned NHPP likelihood is exact]\label{prop:ch06:binnedpoisson}
Let $N$ be an NHPP with deterministic, locally integrable intensity
$\lam(t;\thetamodel)$ on $(0,\Tend]$, and let $I_1,\dots,I_J$ partition
$(0,\Tend]$ with $\Lam_j(\thetamodel) = \int_{I_j}\lam(u;\thetamodel)\dd u$.
Then the bin counts $C_1,\dots,C_J$ are independent with
$C_j \sim \mathrm{Poisson}(\Lam_j)$, and the log-likelihood of the
interval-censored data,
\begin{equation}\label{eq:ch06:binnedpoisson}
\loglik_{\mathrm{IC}}(\thetamodel)
 \;=\;
 \sum_{j=1}^{J}\Big( C_j \log \Lam_j(\thetamodel) - \Lam_j(\thetamodel)
 - \log C_j! \Big),
\end{equation}
is the \emph{exact} likelihood of what was observed --- not an
approximation to the event-time likelihood but the correct likelihood of
the coarsened data. Moreover, if $\lam(\cdot;\thetamodel)$ is constant on
each bin (the weekly piecewise-constant models of \cref{ch:poisson}), the
counts are sufficient for $\thetamodel$ and coarsening loses nothing; in
general the loss is confined to the within-bin shape of $\lam$.
\end{proposition}

\begin{proof}
Start from the density of the ordered event times under the full model: for
$n$ events at $0 < t_1 < \cdots < t_n \le \Tend$,
$f(t_{1:n}) = \prod_{m=1}^{n}\lam(t_m)\,
e^{-\int_0^{\Tend}\lam(u)\dd u}$. The probability of the count vector
$C_{1:J}$ (with $n = \sum_j C_j$) is the integral of $f$ over the
configuration region
$\mathcal{A}(C_{1:J}) = \{\,0<t_1<\cdots<t_n\le\Tend :
\#\{m: t_m \in I_j\} = C_j \ \forall j\,\}$. Because $\lam$ is
deterministic, the integrand is a product of factors each depending on one
coordinate only, so the ordered integral factorizes across bins, and within
bin $j$ the ordered integral of $\prod\lam$ over the $C_j$-dimensional
ordered simplex equals $\Lam_j^{C_j}/C_j!$. Hence
$\Prob(C_{1:J}) = \prod_j \big(\Lam_j^{C_j}/C_j!\big)\,e^{-\Lam_j}$, which
is \cref{eq:ch06:binnedpoisson} and simultaneously the independence and
Poisson-marginal claims. Sufficiency under piecewise-constant $\lam$ is the
factorization criterion applied to \cref{eq:ch06:binnedpoisson}: the
likelihood depends on the data only through $C_{1:J}$.
\end{proof}

\begin{corollary}[For NHPP calibration, untrusted dates cost only resolution]\label{cor:ch06:resolutiononly}
Under \cref{asm:ch06:jitter,asm:ch06:complete}, fitting any NHPP of
\cref{ch:poisson} on bin counts at resolution $\horizon \gg \delta$ uses an
exact likelihood (\cref{prop:ch06:binnedpoisson}) applied to data whose
corruption is bounded by \cref{prop:ch06:jitterrobust}(iii). No
approximation enters anywhere; the entire price of (C4) is the inability to
resolve $\lam$ within bins --- a price we were going to pay anyway by
modeling weekly.
\end{corollary}

This is the sense in which the weekly sector GLM of the current
two-layer architecture (\modrepo{sector\_stability.py}) is already fully
(C4)-compliant: it never consumed positions in the first place. The
classical ancestor of likelihoods under grouped and censored observation is
\citet{turnbull1976}; what is distinctive in our setting is not grouping
per se but the interaction of grouping with \emph{dependence}, which is the
next section's subject.

\section{Interval censoring for Hawkes: where exactness fails}\label{ch06:sec:ichawkes}

Run the proof of \cref{prop:ch06:binnedpoisson} again with the Hawkes
intensity. The density of the ordered times is still the Jacod product, so
the exact interval-censored likelihood is still a configuration integral:
\begin{equation}\label{eq:ch06:icexact}
\lik_{\mathrm{IC}}(\thetamodel)
 \;=\;
 \int_{\mathcal{A}(C_{1:J})}
 \prod_{m=1}^{n} \lam\big(t_m \st \Hist_{t_m^-};\thetamodel\big)\;
 \exp\!\Big({-}\!\int_0^{\Tend}\!\lam\big(u \st \Hist_{u^-};
 \thetamodel\big)\dd u\Big)\;
 \dd t_1 \cdots \dd t_n,
\end{equation}
with $\lam(t\st\Hist_{t^-}) = \mub(t) + \sum_{t_m < t}\kernel(t-t_m)$ and
$n = \sum_j C_j$. But now the factorization step fails, for the precise
reason the model was chosen: the intensity at $t_m$ depends on the
\emph{positions} of all earlier events, including those in earlier bins and
those in the same bin. The integrand does not separate across bins, the
within-bin integrals are not simplex moments of a fixed function, and
\cref{eq:ch06:icexact} is an $n$-dimensional integral --- at the scale of
this dataset, dimension $\sim 5\times 10^4$ --- over a
combinatorially-constrained region, with no product structure to exploit.
Data augmentation (imputing positions by MCMC or EM over
\cref{eq:ch06:icexact}) is possible in principle and is how the Bayesian
literature treats small problems \citep{rasmussen2013}; at our scale, and
for a likelihood that must be refit across sectors, splits, and validation
campaigns, it is not a workable primary route. The likelihood requires
event positions inside bins \emph{because intensity feeds back on
history}; interval censoring withholds exactly that. Something must give,
and the honest thing to give up is the randomness of the intensity ---
replacing it by its mean --- rather than the honesty of the resolution.

\section{The mean-intensity route}\label{ch06:sec:mbpp}

This section is the chapter's centerpiece: the \emph{mean behavior Poisson
process} (MBPP) of \citet{rizoiu2022}, implemented in
\modrepo{mbpp/}. The idea is to fit, instead of the intractable
interval-censored Hawkes likelihood, the exact interval-censored likelihood
of the NHPP whose deterministic intensity is the \emph{expected} Hawkes
intensity. The parameters $(\mub,\branch,\decay)$ are shared, the first
moments are matched exactly, and the price --- discarded dependence and
overdispersion --- is stated and bounded.

\subsection{The mean intensity and its Volterra equation}\label{ch06:sec:volterra}

\begin{theorem}[Renewal equation for the mean intensity]\label{thm:ch06:volterra}
Let $N$ be a Hawkes process on $[0,\Tend]$ with intensity
$\lam(t) = \mub(t) + \int_{(0,t)}\kernel(t-u)\dd N(u)$, where $\mub$ is
deterministic, locally bounded, and $\kernel \ge 0$ with
$\int_0^\infty \kernel = \branch < 1$. Then the mean intensity
$\lambar(t) \defeq \E[\lam(t)]$ is finite and satisfies the linear Volterra
equation of the second kind
\begin{equation}\label{eq:ch06:volterra}
\lambar(t) \;=\; \mub(t) \;+\; \int_0^{t} \kernel(t-u)\,\lambar(u)\dd u .
\end{equation}
\end{theorem}

\begin{proof}
\emph{Step 1 (finiteness).} By the cluster representation of
\cref{ch:hawkes} \citep{hawkesoakes1974}, $N$ is the superposition of
immigrant events arriving with rate $\mub$ and, per event, an independent
Galton--Watson cascade with offspring mean $\branch < 1$; the expected
total progeny per immigrant is $1/(1-\branch)$, so
$\E N(t) \le \int_0^t \mub(u)\dd u \,/\, (1-\branch) < \infty$, and
$\lambar(t) \le \mub(t) + \branch\decay \,\E N(t) < \infty$ for the
exponential kernel (and similarly for any bounded kernel).

\emph{Step 2 (martingale property of the compensator).} For any
\emph{nonnegative predictable} integrand $H$, the defining property of the
intensity gives the compensator formula
$\E\int_0^t H(u)\dd N(u) = \E \int_0^t H(u)\lam(u)\dd u$, both sides
possibly infinite \citep[II.4]{abgk1993}; this is the statement
$\E[\dd N(u)] = \lambar(u)\dd u$ made rigorous. Fix $t$ and apply it to the
deterministic (hence predictable) integrand $H(u) = \kernel(t-u)\ind{u<t}$:
\[
\E\!\int_{(0,t)} \kernel(t-u)\dd N(u)
 \;=\; \E\!\int_0^{t} \kernel(t-u)\,\lam(u)\dd u
 \;=\; \int_0^{t} \kernel(t-u)\,\lambar(u)\dd u ,
\]
the last step by Tonelli (nonnegative integrand) and Step 1. Taking
expectations through
$\lam(t) = \mub(t) + \int_{(0,t)}\kernel(t-u)\dd N(u)$ yields
\cref{eq:ch06:volterra}.
\end{proof}

Note what was used: only the compensator formula --- the same martingale
structure that powers every likelihood in this document --- plus
subcriticality for finiteness. No stationarity, no exponential form, no
constant baseline.

\begin{theorem}[Existence, uniqueness, Neumann series]\label{thm:ch06:neumann}
Under the hypotheses of \cref{thm:ch06:volterra},
\cref{eq:ch06:volterra} has exactly one locally bounded solution on
$[0,\Tend]$, namely
\begin{equation}\label{eq:ch06:neumann}
\lambar \;=\; \mub + R * \mub,
\qquad
R \;\defeq\; \sum_{n\ge 1} \kernel^{*n},
\qquad
\int_0^\infty R \;=\; \frac{\branch}{1-\branch},
\end{equation}
where $\kernel^{*n}$ is the $n$-fold convolution power. Picard iteration
$\lambar_{k+1} = \mub + \kernel * \lambar_k$, $\lambar_0 = \mub$, converges
to it geometrically:
$\sup_{[0,\Tend]}\abs{\lambar_k - \lambar} \le
\branch^{\,k+1}\,(1-\branch)^{-1} \sup_{[0,\Tend]}\mub$.
\end{theorem}

\begin{proof}[Proof sketch]
The map $\Phi(f) = \mub + \kernel * f$ on
$(L^\infty[0,\Tend], \norm{\cdot}_\infty)$ satisfies
$\norm{\Phi(f)-\Phi(g)}_\infty \le \norm{\kernel}_{L^1}
\norm{f-g}_\infty = \branch\norm{f-g}_\infty$, a contraction since
$\branch<1$; Banach's fixed point theorem gives existence, uniqueness among
bounded solutions, and the geometric rate. Unrolling the iteration gives
the Neumann series \cref{eq:ch06:neumann}, whose mass identity follows from
$\norm{\kernel^{*n}}_{L^1} = \branch^n$ and geometric summation. The
general theory of Volterra equations --- including the resolvent
formalism, $L^1$ well-posedness, and what survives when
$\branch \ge 1$ locally --- is \citet{gripenberg1990}. When
$\branch \uparrow 1$ the geometric rate degrades as $(1-\branch)^{-1}$ and
at criticality the series diverges: the mean-intensity route inherits
exactly the stability boundary of the process itself.
\end{proof}

The resolvent $R$ has a probabilistic reading that connects back to
\cref{ch:hawkes}: $\kernel^{*n}$ is the arrival density of $n$-th
generation offspring, so \cref{eq:ch06:neumann} says the mean intensity is
the immigrant rate plus the expected arrival rate of descendants of all
generations --- the cluster representation, averaged.

\subsection{Exponential kernel: closed form}\label{ch06:sec:closedform}

\begin{proposition}[Closed form for the exponential kernel]\label{prop:ch06:expclosed}
Let $\kernel(t) = \branch\decay e^{-\decay t}$ and $\mub(t)\equiv\mub$
constant, with $\branch\in(0,1)$, $\decay>0$. Then
\begin{equation}\label{eq:ch06:expclosed}
\lambar(t)
 \;=\;
 \mub \;+\; \frac{\branch\decay\,\mub}{(1-\branch)\decay}
 \Big(1 - e^{-(1-\branch)\decay\, t}\Big)
 \;=\;
 \mub \;+\; \frac{\branch\,\mub}{1-\branch}
 \Big(1 - e^{-(1-\branch)\decay\, t}\Big),
\end{equation}
so $\lambar$ relaxes from $\lambar(0)=\mub$ to the stationary level
$\lambar(\infty) = \mub/(1-\branch)$ at the \emph{effective relaxation
rate} $r = (1-\branch)\decay$. Equivalently, the resolvent is itself
exponential, $R(t) = \branch\decay\, e^{-(1-\branch)\decay t}$, and for
\emph{time-varying} $\mub(\cdot)$ the solution is the single convolution
$\lambar = \mub + R*\mub$; the integrated version
$\bar\Lam(t) = \int_0^t \lambar$ obeys the same equation,
$\bar\Lam = \int_0^\cdot \mub + \kernel * \bar\Lam$.
\end{proposition}

\begin{proof}
Write $\lambar(t) = \mub + g(t)$ with
$g(t) = \int_0^t \branch\decay\, e^{-\decay(t-u)}\lambar(u)\dd u$. The
kernel's exponential form makes $g$ differentiable with
\[
g'(t) \;=\; \branch\decay\,\lambar(t) \;-\; \decay\, g(t)
 \;=\; \branch\decay\big(\mub + g(t)\big) - \decay\, g(t)
 \;=\; \branch\decay\,\mub \;-\; (1-\branch)\decay\, g(t),
\]
a linear ODE with $g(0)=0$, solved by
$g(t) = \dfrac{\branch\decay\,\mub}{(1-\branch)\decay}
\big(1 - e^{-(1-\branch)\decay t}\big)$; adding $\mub$ gives
\cref{eq:ch06:expclosed}. \emph{Sanity checks.} At $t=0$ the excitation has
not started: $\lambar(0)=\mub$. As $t\to\infty$,
$\lambar \to \mub + \branch\mub/(1-\branch) = \mub/(1-\branch)$ --- the
immigrant rate times the mean cluster size, agreeing with the cluster
representation and with the stationary Hawkes mean of \cref{ch:hawkes};
and substituting the constant $\mub/(1-\branch)$ into
\cref{eq:ch06:volterra} verifies it is the fixed point. For the resolvent
claim, either solve $R = \kernel + \kernel * R$ by the same ODE argument or
sum the Erlang series
$\kernel^{*n}(t) = \branch^n \decay^n t^{n-1}e^{-\decay t}/(n-1)!$
directly; then $\lambar = \mub + R*\mub$ is
\cref{eq:ch06:neumann}. The equation for $\bar\Lam$ follows by integrating
\cref{eq:ch06:volterra} and Fubini.
\end{proof}

The repository implements exactly this structure
(\modrepo{mbpp/core.py}): the impulse response
$R(t) = \branch\decay e^{(\branch-1)\decay t}$ is Eq.~(14) of
\citet{rizoiu2022}; their Theorems~1--2 and Corollary~3 give
$\lambar = s + R * s$ and $\bar\Lam = S + R * S$ in closed form for every
exogenous input $s(\cdot)$ in \modrepo{mbpp/exogenous.py} (constant,
rectangle, piecewise-constant, sinusoidal, impulse trains, and
covariate-driven baselines), with a numerical finite-convolution solver for
kernels without closed forms (power-law). The relaxation rate
$r = (1-\branch)\decay$ of \cref{prop:ch06:expclosed} will reappear twice:
as the decay rate of cross-bin count correlation, and as the parameter
whose product with the bin width governs the $\branch$--$\decay$ ridge.

\subsection{The MBPP working likelihood, and precisely what it discards}\label{ch06:sec:mbpplik}

\begin{definition}[Mean behavior Poisson process and IC likelihood]\label{def:ch06:mbpp}
The MBPP associated with a Hawkes model $(\mub,\kernel)$ is the NHPP with
deterministic intensity $\lambar(\cdot)$ solving
\cref{eq:ch06:volterra}. Its interval-censored log-likelihood on the
partition of \cref{prop:ch06:jitterrobust} is the exact Poisson form
\cref{eq:ch06:binnedpoisson} with means
$\bar\Lam_j(\thetamodel) = \int_{I_j}\lambar(u;\thetamodel)\dd u$:
\begin{equation}\label{eq:ch06:mbppll}
\loglik_{\mathrm{MBPP}}(\thetamodel)
 \;=\;
 \sum_{j=1}^{J} \Big( C_j \log \bar\Lam_j(\thetamodel)
 \;-\; \bar\Lam_j(\thetamodel) \Big) \;+\; \text{const},
\end{equation}
maximized over $\thetamodel = (\mub, \branch, \decay, \dots)$
(\modrepo{mbpp/interval\_censored.py}, \code{fit\_mbpp\_ic}).
\end{definition}

We state exactly what \cref{eq:ch06:mbppll} gets right and wrong when the
data is generated by the Hawkes process itself.

\emph{Exact: the first moments.} By \cref{thm:ch06:volterra} and the
compensator formula, $\E C_j = \E\int_{I_j}\lam = \bar\Lam_j$
\emph{exactly}, for every $j$, transient or stationary. Consequently the
expected score of \cref{eq:ch06:mbppll} at the true parameters is
$\sum_j (\E C_j/\bar\Lam_j - 1)\partial_{\thetamodel}\bar\Lam_j = 0$: the
MBPP maximum-likelihood estimator is an M-estimator with an unbiased
estimating equation --- a moment-matching estimator dressed as a
likelihood --- and is consistent under the usual regularity and ergodicity
conditions for correctly specified means. \citet[Prop.~8]{rizoiu2022} give
the population version: the expected IC log-likelihood is (up to constants)
the generalized Kullback--Leibler divergence between the true interval
means and the model's interval compensators, minimized exactly at
parameters reproducing the true means --- identifiability is preserved
\emph{in the KL sense}, to the extent that the interval means themselves
separate the parameters (they pin $\mub$ and $\branch$ through levels and
transients; how much of $\decay$ they retain is
\cref{ch06:sec:ridge}).

\emph{Discarded: the second moments.} Two falsehoods are asserted by
\cref{eq:ch06:mbppll}. First, that $C_j$ is Poisson --- but the Hawkes
random compensator $\int_{I_j}\lam$ has been replaced by its expectation,
and the law of total variance charges for it:
$\Var(C_j) = \E[\text{(martingale variation)}] + \Var\big(\textstyle
\int_{I_j}\lam\big) + \text{cross terms} > \E C_j$. Second, that
$C_1,\dots,C_J$ are independent --- but excitation propagates across bin
boundaries, and for the exponential kernel the count autocovariance decays
geometrically with lag at exactly the AR root $e^{-r\horizon}$,
$r = (1-\branch)\decay$. The next proposition quantifies the first
falsehood; the second is smallest precisely when bins are wide relative to
the relaxation time $1/r$.

\begin{proposition}[Hawkes bin counts are overdispersed]\label{prop:ch06:overdisp}
For a stationary Hawkes process with $\kernel = \branch\decay
e^{-\decay t}$, $\branch\in(0,1)$, constant $\mub$, let
$F(\horizon) = \Var(C_j)/\E(C_j)$ be the Fano factor of a bin of width
$\horizon$. Then
\begin{equation}\label{eq:ch06:fano}
F(\horizon) \;\longrightarrow\; 1 \quad (\horizon \to 0),
\qquad
F(\horizon) \;\longrightarrow\; \frac{1}{(1-\branch)^2}
\quad (\horizon \to \infty),
\end{equation}
and $F(\horizon) > 1$ for all $\horizon > 0$: true bin counts are strictly
overdispersed relative to the Poisson law asserted by
\cref{eq:ch06:mbppll}, by a factor approaching $(1-\branch)^{-2}$ for wide
bins.
\end{proposition}

\begin{proof}[Proof sketch]
\emph{Wide bins.} By the cluster representation
\citep{hawkesoakes1974}, $N$ is a superposition of i.i.d.\ clusters
attached to Poisson($\mub$) immigrants; the total cluster size $S$ is the
total progeny of a Galton--Watson process with Poisson($\branch$)
offspring, for which $\E S = (1-\branch)^{-1}$ and
$\Var S = \branch(1-\branch)^{-3}$. Cluster durations have finite mean
(exponential delays), so as $\horizon\to\infty$ the count over a bin is,
up to boundary terms of $O(1)$ against a mean of $O(\horizon)$, a compound
Poisson sum $\sum_{i\le K} S_i$ with $K\sim$ Poisson($\mub\horizon$);
hence
$F \to \E S^2/\E S = \big(\Var S + (\E S)^2\big)/\E S =
\branch(1-\branch)^{-2} + (1-\branch)^{-1} = (1-\branch)^{-2}$.
\emph{Narrow bins.} $\E[C_j(C_j-1)] = \iint_{I_j^2}\rho_2 = O(\horizon^2)$
for the (bounded) second-order product density $\rho_2$ of the stationary
process, while $\E C_j = \Theta(\horizon)$, so
$F = 1 + O(\horizon) \to 1$. Strict inequality for all $\horizon$ follows
from positivity of the covariance density of a self-exciting process with
nonnegative kernel \citep{hawkes1971,daley2003}. What breaks if
$\branch \ge 1$: $\E S = \infty$ and the Fano factor diverges --- another
face of the stability boundary.
\end{proof}

\begin{warning}[MBPP standard errors are optimistic]\label{warn:ch06:optimisticse}
The observed-information standard errors of \cref{eq:ch06:mbppll} (as
reported by \code{fit\_mbpp\_ic}) price the counts as Poisson and
independent. By \cref{prop:ch06:overdisp} the true variance is up to
$(1-\branch)^{-2}$ larger per bin --- at $\branch = 0.35$, for example,
this derived ceiling is about $2.4$, so na\"ive standard errors can be
about $1.5\times$ too small --- and positive cross-bin correlation
compounds the understatement. Remedies, in committed order: (i) sandwich
(robust M-estimation) covariance, $\Fisher^{-1} V \Fisher^{-1}$ with $V$
the empirical score variance --- valid because the estimating equation is
unbiased regardless of the working variance \citep[Ch.~5]{vandervaart1998};
(ii) a negative-binomial working likelihood with variance
$\bar\Lam_j + \bar\Lam_j^2/\nu$ when the dispersion diagnostic of
\cref{ch:gof} (Pearson statistic, \modrepo{mbpp/gof.py}) exceeds the
predicted Fano band; (iii) the Bayesian route of \modrepo{mbpp/bayes.py},
whose posterior exhibits the $\branch$--$\decay$ ridge as a correlation
instead of hiding it in a scalar. Never report the na\"ive SEs alone.
\end{warning}

\emph{When is the mean-intensity replacement good?} Three regimes, all of
which hold for our intended use. (i) \emph{Aggregation over many
low-feedback units:} sector-level counts superpose many firms' sparse
streams; superpositions of many sparse, weakly dependent point processes
are close to Poisson \citep[Ch.~9]{daley2003}, so the Poisson working law
is close to true even before the mean replacement. (ii) \emph{Subcritical,
moderate $\branch$:} the entire second-moment error is bounded by the Fano
ceiling $(1-\branch)^{-2}$, which is $O(1)$ and known; the bound stays
modest as long as the fitted $\branch$ sits well below one, and the fitted
values will be reported in \cref{ch:results-spain}.
(iii) \emph{Bin width at or beyond the kernel half-life:} the neglected
cross-bin dependence has AR root $e^{-r\horizon}$, small once
$\horizon \gtrsim 1/r$; and within-bin position information --- the thing
MBPP cannot use --- was already destroyed by (C4). The regime where MBPP is
\emph{bad} is the mirror image: a single high-$\branch$ stream at fine
resolution, where one should be fitting the event-time likelihood of
\cref{ch:hawkes} instead --- on verified dates.

\subsection{What survives aggregation: $\branch$ yes, $\decay$ no}\label{ch06:sec:ridge}

The stationary interval mean is
$\bar\Lam_j = \mub\horizon/(1-\branch)$: it contains $\branch$ and
\emph{no trace of} $\decay$. All $\decay$-information in counts therefore
lives in temporal structure --- transients of
\cref{eq:ch06:expclosed} and cross-bin autocovariances --- and both enter
only through the relaxation root $e^{-r\horizon}$. A Fisher-information
computation makes this quantitative: the information for $\decay$ carried
by counts decays like $\horizon^{2}e^{-2(1-\branch)\decay\horizon}$ ---
exponentially in the bin width --- so binning does not merely blur the
timescale; past a few half-lives it erases it. This yields a working rule,
stated here as a hypothesis and scheduled for validation on synthetic data
in \cref{ch:results-synth}: \emph{bucket width $\le$ one kernel half-life
should be safe for $\decay$; $\ge$ four half-lives should destroy timing;
$\branch$ and covariate effects, read from first moments, should survive
any width}. The expected failure signature at coarse widths is a $\decay$
biased toward the ridge with a deceptively small spread --- the mark of an
unidentified direction --- and the planned recovery experiments must
reproduce it before the rule is trusted. The Bayesian implementation
(\modrepo{mbpp/bayes.py}) is the honest reporting vehicle on the ridge: a
flat-prior posterior on $(\branch,\decay)$ stays wide along $\decay$ and
tight along $\branch$, and a $\decay$ prior from a verified-date subset
tightens it legitimately.

\begin{warning}[Never claim timing precision finer than the honest resolution]\label{warn:ch06:timingclaims}
A $\decay$ fitted from weekly counts is a ridge coordinate, not a
measurement. Design consequence, binding on every downstream consumer
(\cref{ch:decision}): no model output may assert within-week timing
structure --- expected response times, decay half-lives, ``momentum
windows'' --- unless $\decay$ was estimated from a verified-date subset at
event-time resolution, and the honest resolution of the input data is
reported next to any timing claim. The failure mode is tempting because
the fit \emph{converges}: the optimizer returns a $\decay$ with a small
na\"ive SE while sliding along the ridge.
\end{warning}

\section{Single untrusted dates as one-bin censoring}\label{ch06:sec:onebin}

The machinery above treats wholesale aggregation. The same logic handles
retail corruption: a \emph{single} deal whose date is trusted only to
$\pm\delta_m$ is an event observed in the one-off interval
$I_m = (\hat t_m - \delta_m,\, \hat t_m + \delta_m]$, i.e.\ a private
one-bin censoring. Mixed records --- verified dates for some deals,
windows for others --- then call for a mixed likelihood.

\begin{proposition}[Exact mixed likelihood, Poisson case]\label{prop:ch06:mixedpoisson}
Let $N$ be an NHPP with deterministic intensity $\lam$, and suppose the
record consists of (a) verified events at exact times
$\{t_m\}_{m\in\mathcal{V}}$ and (b) untrusted events known only to lie in
windows $\{I_m\}_{m\in\mathcal{U}}$, one event per window, where the
windows are pairwise disjoint and contain no verified events (overlapping
windows are first merged into a single interval carrying its total count,
reducing to \cref{prop:ch06:binnedpoisson} locally), and (c) the assertion,
by \cref{asm:ch06:complete}, that no further events occurred in
$(0,\Tend]$. The exact log-likelihood of this record is
\begin{equation}\label{eq:ch06:mixed}
\loglik_{\mathrm{mix}}(\thetamodel)
 \;=\;
 \sum_{m\in\mathcal{V}} \log\lam(t_m;\thetamodel)
 \;+\;
 \sum_{m\in\mathcal{U}} \log \Lam(I_m;\thetamodel)
 \;-\;
 \int_0^{\Tend}\lam(u;\thetamodel)\dd u,
\qquad
\Lam(I) = \int_I \lam .
\end{equation}
\end{proposition}

\begin{proof}
The restrictions of an NHPP to disjoint sets are independent. Partition
$(0,\Tend]$ into the windows $\{I_m\}_{m\in\mathcal U}$ and their
complement $B$. On $B$ the observation is exact: points at
$\{t_m\}_{m\in\mathcal V}$ and none elsewhere, with density
$\prod_{\mathcal V}\lam(t_m)\, e^{-\Lam(B)}$. On each $I_m$ the
observation is the event $\{N(I_m)=1\}$, with probability
$\Lam(I_m)e^{-\Lam(I_m)}$. Multiplying and collecting exponents (the
$\Lam$'s add to $\int_0^{\Tend}\lam$ over the partition) gives
\cref{eq:ch06:mixed}. Every factor is exact; no within-window position was
used.
\end{proof}

So for Poisson calibration, mixed trust costs nothing but the windows'
internal resolution: verified deals keep their full event terms, untrusted
deals degrade gracefully to interval terms, and the survival factor is
common. For Hawkes the by-now-familiar obstacle returns --- the intensity
after an untrusted event depends on where in $I_m$ it actually happened ---
and the exact mixed likelihood is again a configuration integral over
$\prod_{m\in\mathcal U} I_m$. The committed practical recipe, flagged as
an approximation, mirrors the MBPP construction locally: keep exact terms
where dates are verified, and push untrusted events through their
\emph{expected} kernel contribution. Concretely, define the smeared
intensity
\begin{equation}\label{eq:ch06:smeared}
\tilde\lam(t;\thetamodel)
 \;=\;
 \mub(t)
 \;+\; \sum_{m\in\mathcal{V}:\, t_m<t} \kernel(t - t_m)
 \;+\; \sum_{m\in\mathcal{U}} \frac{1}{\abs{I_m}}
 \int_{I_m} \kernel(t-u)\,\ind{u<t}\dd u,
\end{equation}
(the untrusted event's excitation averaged uniformly over its window, with
only the past part of the window contributing), and maximize \cref{eq:ch06:mixed} with
$\tilde\lam$ in place of $\lam$. This is exact as $\branch\to0$
(Poisson limit), errs only in the untrusted events' second-order feedback
terms, and must be validated by the protocol of \cref{ch06:sec:validation}
before use --- the same discipline as the MBPP itself. When untrusted
events are dense in some sector-period, the recipe degenerates gracefully:
merge that period's windows into bins and the untrusted sum in
\cref{eq:ch06:mixed} becomes the MBPP likelihood \cref{eq:ch06:mbppll}
restricted to that period.

\section{The validation loop}\label{ch06:sec:validation}

Every approximation admitted in this chapter --- the mean-intensity
replacement, the mixed-likelihood smearing, the choice of resolution ---
is validated by the same closed loop, which the repository already
implements for the MBPP (\modrepo{mbpp/ic\_simulate.py}) and which is
hereby made the standing acceptance test.

\begin{protocol}[Simulate exactly, censor, refit, compare]\label{prot:ch06:icloop}
\leavevmode
\begin{enumerate}[label=(\alph*)]
\item \emph{Simulate exactly} from the continuous-time Hawkes model at the
fitted (or stress) parameters: the cluster/branching construction of
\citet{hawkesoakes1974} (\code{simulate\_separable\_hawkes}), or exact
simulation by \citet{dassioszhao2013}, or Ogata thinning
\citep{ogata1981,lewisshedler1979} --- never from the MBPP itself, which
would assume the conclusion.
\item \emph{Censor as production does}: collapse event times to the
production partition with \code{interval\_censor}; optionally inject
adversarial jitter at the audited $\delta$ before binning, to certify
\cref{prop:ch06:jitterrobust} empirically.
\item \emph{Refit with the production likelihood}
(\code{fit\_mbpp\_ic} and variants; the mixed recipe of
\cref{ch06:sec:onebin} where applicable).
\item \emph{Compare}: (1) parameter recovery --- $\branch$ must be
recovered within Monte Carlo bands at every candidate resolution;
$\decay$ only at resolutions the ridge rule permits, and its failure at
coarse bins should reproduce the ridge signature predicted in
\cref{ch06:sec:ridge}; (2) SE honesty --- the ratio of Monte Carlo spread to
reported SE is the empirical dispersion factor, to be compared against the
Fano prediction of \cref{prop:ch06:overdisp} and repaired per
\cref{warn:ch06:optimisticse}; (3) fit diagnostics --- Pearson dispersion
and PIT-style count calibration \citep{czado2009} via
\modrepo{mbpp/gof.py}, with the event-time time-rescaling check
\citep{brown2002} reserved for verified-date subsets
(\cref{ch:gof}).
\item \emph{Gate}: a resolution/likelihood pair is admitted for real-data
use only after passing (d) in the synthetic world at matched sample sizes
--- the transfer convention of \cref{ch02:sec:synthetic}. Campaign
results will be reported in \cref{ch:results-synth}.
\end{enumerate}
\end{protocol}

\section{Defect declaration and commitments}\label{ch06:sec:decisions}

Per \cref{tab:ch02:defects} and decision D2.1, this chapter \emph{owns
(C4)} and relates to the rest of the ledger as follows.
\emph{(C1) corrected here}: \cref{thm:ch06:rightcensor,prop:ch06:locality}
establish that right censoring costs nothing for history-dependent
intensities when open spells are kept, closing the Hawkes case left open by
\cref{warn:ch02:openspells}. \emph{(C2) respected mechanically}: every
integral and every bin in this chapter starts at the entry time
$\entry_i$ where firm-level objects are involved; the selection content of
(C2) is owned by \cref{ch:universe}. \emph{(C3) vulnerable, out of
scope}: nothing here repairs contaminated risk sets --- sector-level count
likelihoods are largely insulated (sectors do not die), but any firm-level
use of these tools inherits (C3) with the deflation bias documented in
\cref{ch:data}, owned by \cref{ch:poisson,ch:buckets,ch:ranking}.
\emph{(C5) orthogonal}: covariates enter the MBPP through the exogenous
input $s(\cdot)$ (\modrepo{mbpp/exogenous.py}); their staleness is
\cref{ch:stale}'s problem and its remedies compose with everything here.
\emph{(C6) robust at sector level}: sector deal counts include
out-of-universe deals, so sector-level MBPP consumes complete counts;
firm-attributable applications inherit (C6) as thinning, in tension with
\cref{asm:ch06:complete}, and are owned by \cref{ch:universe}.

\begin{decision}{D6.1 --- the honest-resolution rule}
The observation resolution for continuous-time ground-process modeling on
PitchBook 2011--2024 is \emph{weekly}. Daily resolution is permitted only
for subsets whose closed dates the date audit certifies with
$\delta \le 1$ day; back-filled and announced-only records never qualify.
Rationale: \cref{prop:ch06:jitterrobust} (boundary-transfer error
$\sim 4\delta/\horizon$), \cref{asm:ch06:convention} (announced-vs-closed
spreads of days, heaping in backfills), and the ridge rule of
\cref{ch06:sec:ridge}: for kernel half-lives near one week, weekly bins
should keep $\branch$ identifiable and leave $\decay$ marginal, a
hypothesis \cref{ch:results-synth} will test. Reversal trigger: a
data vintage whose audit certifies $\delta \le 1$ day for a sector's
majority of deals moves that sector to daily; conversely, evidence of
heaping beyond month boundaries (quarter-level backfill) forces that
segment to monthly.
\end{decision}

\begin{decision}{D6.2 --- MBPP is the committed aggregation likelihood}
Any sector-level continuous-time Hawkes model fitted on binned counts uses
the MBPP interval-censored likelihood \cref{eq:ch06:mbppll}
(\code{fit\_mbpp\_ic} and variants), with standard errors repaired per
\cref{warn:ch06:optimisticse} (sandwich by default; negative-binomial or
Bayesian escalation on a failed dispersion check), and with every fit
gated by \cref{prot:ch06:icloop}. Rationale: exact first moments
(\cref{thm:ch06:volterra}), KL-sense identifiability
\citep{rizoiu2022}, closed forms for the exponential kernel
(\cref{prop:ch06:expclosed}), and a bounded, quantified price
(\cref{prop:ch06:overdisp}) --- against an exact alternative
\cref{eq:ch06:icexact} that is computationally out of reach. Reversal
trigger: if dispersion diagnostics persistently exceed the Fano ceiling
$(1-\hat\branch)^{-2}$ (pointing at latent-frailty variance the MBPP
cannot absorb), the sector-level count model reverts to the discrete
autoregressive count family of \cref{ch:buckets}
\citep{fokianos2009}, which models the conditional variance directly.
\end{decision}

\begin{decision}{D6.3 --- event-time likelihoods reserved for verified dates}
The exact event-time likelihood of \cref{ch:hawkes} (and the time-rescaling
diagnostics that need it) may be fitted only on date-audit-verified
subsets; $\decay$ estimated anywhere else is prior-constrained from such a
subset or fixed (decay banks). Mixed records use
\cref{eq:ch06:mixed}: exactly for Poisson
(\cref{prop:ch06:mixedpoisson}), with the flagged smeared-intensity recipe
\cref{eq:ch06:smeared} for Hawkes. Rationale:
\cref{warn:ch06:timingclaims} --- timing claims must trace to data that
contains timing. Reversal trigger: none for the principle; the boundary of
``verified'' moves with the date audit.
\end{decision}

\begin{codehooks}
The mathematics of this chapter maps onto \modrepo{mbpp/} as follows.
\emph{\code{core.py}}: the MBPP object --- kernels, the resolvent/impulse
response $R(t)=\branch\decay e^{(\branch-1)\decay t}$ of
\cref{prop:ch06:expclosed}, closed-form and numeric solutions of
\cref{eq:ch06:volterra}, interval compensators $\bar\Lam_j$, and the
Proposition-6/7 compensator bounds used for forecasting.
\emph{\code{interval\_censored.py}}: the IC likelihood
\cref{eq:ch06:mbppll} (\code{ic\_ll}), the SSE/Bregman alternative, and
the fitting family \code{fit\_mbpp\_ic} /
\code{\_multi} / \code{\_covariates} / \code{\_sumexp} /
\code{\_excitation}, plus \code{forecast\_counts}.
\emph{\code{exogenous.py}}: the exogenous inputs $s(\cdot)$ ---
piecewise-constant, covariate-driven, impulse, LHPP --- through which
\cref{ch:stale}'s covariates and \cref{ch:poisson}'s baselines enter.
\emph{\code{ic\_simulate.py}}: the branching simulator and
\code{interval\_censor}, powering \cref{prot:ch06:icloop}.
\emph{\code{gof.py}}: Pearson/dispersion diagnostics for counts and
time-rescaling residuals for verified-date subsets;
\emph{\code{bayes.py}}: the posterior treatment of the
$\branch$--$\decay$ ridge. \emph{To be written}:
\modrepo{data/date\_audit.py} --- the audit behind
\cref{asm:ch06:jitter,asm:ch06:convention}: classify each deal date as
closed/announced/backfilled, estimate convention offsets $b(r)$, detect
month/quarter heaping, emit per-event trust radii $\delta_m$ and the
verified-date flags consumed by D6.1/D6.3 --- and the sandwich-covariance
wrapper for \code{fit\_mbpp\_ic} required by
\cref{warn:ch06:optimisticse}.
\end{codehooks}
